We first remove the continuous orbit optimisation exactly.
For complex diagonal matrices $D_v,D_w$ of any common size,
\begin{equation}\label{eq:fqe-signed-diagonal-overlap}
 \max_{A,B\in O(m)}|\operatorname{tr}(D_v^\dagger A D_wB^T)|
 =\max_{\pi,\,\varepsilon_i=\pm1}
                  \left|\sum_i\bar v_i\varepsilon_iw_{\pi(i)}\right|.
\end{equation}
Indeed $|A_{ij}B_{ij}|$ is doubly substochastic, by Cauchy--Schwarz
on each row and column. Complete it to a doubly stochastic matrix;
for any common phase the real overlap is bounded by the largest
assignment sum of the absolute real coefficient entries. Birkhoff's
permutation decomposition gives this bound, and signed permutation
matrices attain it. Maximising the phase proves
\eqref{eq:fqe-signed-diagonal-overlap}.

The magic normal form makes a critical state $D_v$. Here and below the coordinates of $v\in\mathbb C^4$, and the
rows and columns of $4\times4$ matrices, are indexed $0,\dots,3$: $v_i$ is the coefficient of $m_{i+1}\otimes m_{i+1}$ in
$G(v)$ of Appendix~\ref{sec:fqe-normal-form} of the paper, so $D_v=\operatorname{diag}(v)$.
Up to global phase, the union $\mathcal M$ is the $O(4)\times O(4)$ orbit
of $Z_*$: a common orientation reversal in both factors fixes $Z_*$
up to proper rotations, whereas a single reversal gives its conjugate
up to proper rotations. Thus its overlap with $D_v$ is exactly the
maximum over the following 32 diagonal representatives:
\[
 \mathcal H=\left\{\eta:\ \eta_i=1/\sqrt2,\quad
         \eta_j=\varepsilon_j i/\sqrt6\ (j\ne i),\quad i=0,1,2,3\right\}.
\]
Its exact spin-flip-dependent covering profile is proved in
Lemma~\ref{lem:fqe-critical-real-density} below.

The real cut transforms are $H_0=I$ and
\[
 H_1=\frac12\begin{pmatrix}
 1&-1&1&-1\\-1&1&1&-1\\1&1&1&1\\-1&-1&1&1
 \end{pmatrix},\qquad
 H_2=\frac12\begin{pmatrix}
 1&-1&-1&-1\\-1&1&-1&-1\\1&1&-1&1\\-1&-1&-1&1
 \end{pmatrix}.
\]
With $\Delta=\operatorname{diag}(1,-1,1,-1)$, so that $\Delta v$ are the coordinates $(a,b,c,d)$ of $G_{abcd}$, these are
the transforms $H_\pm$ of Section~\ref{sec:sup-sextic-zeros} conjugated by $\Delta$: $H_1=\Delta H_+\Delta$, and $H_2$
is $\Delta H_-\Delta$ with its rows permuted and re-signed, which only permutes the squared moduli $|e_i^TH_2v|^2$ of the cut spectrum.

\begin{proof}[Proof of Theorem~\ref{thm:fqe-critical-orbit-sharp}]
The comparison $Q(\psi)\ge1-\Phi(\psi)$ holds for all states by
Theorem~\ref{thm:fqe-global-overlap}. When $\Phi(\psi)\ge\tfrac12$,
$1-\Phi(\psi)\ge\tfrac32-2\Phi(\psi)$, so \eqref{eq:fqe-critical-fidelity}
holds. Assume henceforth that $\Phi(v)\le\tfrac12$ for a unit critical
normal vector $v\in\mathbb C^4$.

The three real orthogonal cut transforms $H_0=I,H_1,H_2$ displayed above
preserve $N=\|v\|^2=1$, $\tau=|\sum_i v_i^2|^2$ and $Q(v)$, and permute
the $32$ representatives $\mathcal H$ up to the sixth-root-of-unity
phases $(\pm1\pm i\sqrt3)/2$: for any $\eta\in\mathcal H$, one coordinate
is $1/\sqrt2$ and three are $\pm i/\sqrt6$, so under $H_1$ or $H_2$ the
row whose relative signs on those three coordinates match $\eta$ has inner
product $\frac{\pm1\pm i\sqrt3}{2\sqrt2}$ and the other three rows have
$\frac{\pm1\pm i\sqrt3}{2}(\mp i/\sqrt6)$ (CO-1). Thus $H_0,H_1,H_2$,
coordinate permutations, global phase and coordinate sign flips preserve
both $Q(v)$ and $\Phi(v)$. Replace $v$ by $H_mv$ ($m\in\{0,1,2\}$) and
permute coordinates so that
$p_0:=|v_0|^2=\max_{m,i}|e_i^TH_mv|^2$ is the largest of all twelve cut
probabilities. Choose global phase so that $v_0=r_0>0$, and flip signs
$v_j\mapsto\pm v_j$ ($1\le j\le3$) so that $y_j:=\operatorname{Im}v_j\ge0$,
writing $v_j=x_j+iy_j$ and $p_j=x_j^2+y_j^2\le p_0=r_0^2$. Put
$X_2:=\sum_{j=1}^3x_j^2$, $Y_2:=\sum_{j=1}^3y_j^2=1-p_0-X_2$,
$S_x:=\sum_{j=1}^3|x_j|$ and $S_y:=\sum_{j=1}^3y_j$.

Testing $\Phi(v)$ against
$Z_*=(1/\sqrt2,i/\sqrt6,i/\sqrt6,i/\sqrt6)^T\in\mathcal H$ gives
\[
 \Phi(v)\ge\Phi_{\rm im}(v):=\left(\frac{r_0}{\sqrt2}+\frac{S_y}{\sqrt6}\right)^2
 =\frac12-\frac12X_2+\frac13\Bigl(\sqrt3\,r_0S_y-Y_2+\sum_{1\le k<\ell\le3}y_ky_\ell\Bigr).
\]
Since $\Phi_{\rm im}(v)\le\Phi(v)\le\tfrac12$, we have
$S_y\le\sqrt3(1-r_0)$, hence $Y_2\le S_y^2\le3(1-r_0)^2$ and
$X_2=1-r_0^2-Y_2\ge2(1-r_0)(2r_0-1)$. Moreover, since
$0\le y_j\le\sqrt{p_j}\le r_0$ and $\sqrt3<\tfrac{26}{15}<\tfrac74$
($26^2=676>675$), we also have
$Y_2\le r_0S_y\le\tfrac{26}{15}r_0(1-r_0)\le\tfrac74r_0(1-r_0)$ and
$X_2\ge\tfrac1{15}(1-r_0)(15-11r_0)$.

Up to sign, the eight rows of $H_1$ and $H_2$ are
$\tfrac12(1,\varepsilon_1,\varepsilon_2,\varepsilon_3)$ for
$\varepsilon\in\{\pm1\}^3$, so the eight cut probabilities in $H_1v$ and
$H_2v$ are
\[
 \lambda_\varepsilon
 =\frac14\left[\left(r_0+\sum_{j=1}^3\varepsilon_jx_j\right)^2
   +\left(\sum_{j=1}^3\varepsilon_jy_j\right)^2\right]
 =\frac14\left(1+2\sum_{j=1}^3\varepsilon_jc_{0j}
   +2\sum_{1\le k<\ell\le3}\varepsilon_k\varepsilon_\ell c_{k\ell}\right)\le p_0,
\]
where $c_{ij}:=\operatorname{Re}(\bar v_iv_j)=x_ix_j+y_iy_j$
($c_{0j}=r_0x_j$). Choosing $\varepsilon_j=\operatorname{sgn}(x_j)$ gives
$(r_0+S_x)^2\le4p_0$, so $S_x\le r_0$. For each $j\in\{1,2,3\}$, averaging
$\lambda_\varepsilon\le p_0$ over the four sign vectors with
$\varepsilon_j=\operatorname{sgn}(x_j)$ cancels the other terms and yields
$\tfrac14(1+2r_0|x_j|)\le r_0^2$, so $|x_j|\le\frac{4r_0^2-1}{2r_0}$ and
$X_2\le(\max_j|x_j|)S_x\le2r_0^2-\tfrac12$. Combining this with
$Y_2\le\tfrac74r_0(1-r_0)$ gives
\[
 1=r_0^2+X_2+Y_2\le3r_0^2-\frac12+\frac74r_0(1-r_0),\qquad
 0\le\frac14(5r_0-3)(r_0+2),
\]
hence $r_0\ge\tfrac35$. Furthermore, for $\varepsilon_j=\operatorname{sgn}(x_j)$
we have $(\sum_{j=1}^3\varepsilon_jy_j)^2\ge Y_2-2\sum_{k<\ell}y_ky_\ell
=2Y_2-S_y^2\ge2(1-p_0-X_2)-3(1-r_0)^2$, so
\[
 (r_0+S_x)^2\le4p_0-2(1-p_0-X_2)+3(1-r_0)^2=(3r_0-1)^2+2X_2
 \le\left(3r_0-1+\frac{X_2}{3r_0-1}\right)^2,
\]
and therefore $S_x\le\min\bigl(r_0,\,2r_0-1+\frac{X_2}{3r_0-1}\bigr)$. For
$s_2:=\sum_{1\le k<\ell\le3}|x_k||x_\ell|=\tfrac12(S_x^2-X_2)$, if
$X_2\le(1-r_0)(3r_0-1)$ then
\[
 2s_2\le\left(2r_0-1+\frac{X_2}{3r_0-1}\right)^2-X_2
 =(2r_0-1)^2-\frac{X_2\bigl((1-r_0)(3r_0-1)-X_2\bigr)}{(3r_0-1)^2}
 \le(2r_0-1)^2,
\]
while if $X_2\ge(1-r_0)(3r_0-1)$ then
$2s_2\le r_0^2-X_2\le(2r_0-1)^2$. Thus $2s_2\le(2r_0-1)^2$ and
$q_4:=\sum_{1\le k<\ell\le3}x_k^2x_\ell^2\le s_2^2\le\tfrac14(2r_0-1)^4$
for all $X_2\ge0$.

On the critical locus, expanding $Q(v)$ in terms of $p_i=|v_i|^2$ and
$c_{ij}=\operatorname{Re}(\bar v_iv_j)$ gives the exact identity (CO-1)
\begin{equation}\label{eq:fqe-critical-three-term}
 Q(v)=\frac12\Bigl(\sum_{i=0}^3p_i\Bigr)^3
      +18\sum_{j=1}^3D_j\bigl(c_{0j}^2-c_{k\ell}^2\bigr),
\end{equation}
where for each $j\in\{1,2,3\}$ and $\{k,\ell\}=\{1,2,3\}\setminus\{j\}$ we
set $D_j:=p_0+p_j-p_k-p_\ell=2(p_0+p_j)-1=1-2(p_k+p_\ell)\le
D_{\max}:=\min(1,4p_0-1)$. Using $c_{0j}^2=p_0x_j^2$,
$c_{k\ell}^2\le2(x_k^2x_\ell^2+y_k^2y_\ell^2)$ when $D_j\ge0$, and
$c_{k\ell}^2\ge0$ when $D_j<0$, we obtain with $D_j^+:=\max(D_j,0)$ the
decoupled lower bound
\[
 Q(v)+2\Phi(v)-\frac32\ge Q(v)+2\Phi_{\rm im}(v)-\frac32
 \ge\mathcal Y(r_0,y)+\mathcal X(p;x^2),
\]
where
\begin{align*}
 \mathcal Y(r_0,y)&:=\frac23\Bigl(\sqrt3\,r_0S_y-Y_2
   +\sum_{1\le k<\ell\le3}y_ky_\ell\Bigr)-36\sum_{j=1}^3D_j^+y_k^2y_\ell^2,\\
 \mathcal X(p;x^2)&:=\sum_{j=1}^3(18p_0D_j-1)x_j^2
   -36\sum_{j=1}^3D_j^+x_k^2x_\ell^2.
\end{align*}

We show that both $\mathcal Y(r_0,y)$ and $\mathcal X(p;x^2)$ are
nonnegative (CO-2). For $\mathcal Y(r_0,y)$, the four-term AM--GM
inequality on $(2r_0+1)+3(2r_0-1)+4(1-r_0)+4(1-r_0)=6$ gives
$48(4r_0^2-1)(1-r_0)^2\le(3/2)^4=81/16$, so
$9D_{\max}S_y^2\le27(4r_0^2-1)(1-r_0)^2\le729/256<3$. Since
$Y_2=S_y^2-2\sum_{k<\ell}y_ky_\ell$ and
$\sum_{k<\ell}y_k^2y_\ell^2\le\tfrac14S_y^2\sum_{k<\ell}y_ky_\ell$
(as $y_ky_\ell\le\tfrac14(y_k+y_\ell)^2\le\tfrac14S_y^2$), we have
\[
 \mathcal Y(r_0,y)\ge\frac23S_y(\sqrt3\,r_0-S_y)
   +\bigl(2-9D_{\max}S_y^2\bigr)\sum_{1\le k<\ell\le3}y_ky_\ell
 \ge\frac23S_y(\sqrt3\,r_0-S_y)-\sum_{1\le k<\ell\le3}y_ky_\ell.
\]
Using $\sum_{k<\ell}y_ky_\ell\le\tfrac13S_y^2$ and
$S_y\le\sqrt3(1-r_0)$ gives
$\mathcal Y(r_0,y)\ge\tfrac13S_y(2\sqrt3\,r_0-3S_y)
\ge\tfrac{\sqrt3}{3}S_y(5r_0-3)\ge0$.

For $\mathcal X(p;x^2)$, substituting
$D_j=(2p_0-1)+2x_j^2+2y_j^2\ge(2p_0-1)+2x_j^2$ and
$\sum_{j=1}^3x_j^4=X_2^2-2q_4$ into the first sum yields
\[
 \mathcal X(p;x^2)\ge\bigl(18p_0(2p_0-1)-1\bigr)X_2+36p_0(X_2^2-2q_4)
   -36\sum_{j=1}^3D_j^+x_k^2x_\ell^2.
\]
If $X_2=0$, then $\mathcal X=0$; assume $X_2>0$. On
$r_0\in[\tfrac{19}{25},1]$, we have $p_0\ge\tfrac{361}{625}>\tfrac12$, so
$D_j=1-2(p_k+p_\ell)\ge0$ and $D_j\le1-2(x_k^2+x_\ell^2)$. Hence
$\sum_{j=1}^3D_j^+x_k^2x_\ell^2\le q_4-2(X_2q_4-3x_1^2x_2^2x_3^2)
\le(1-\tfrac43X_2)q_4$ by $3x_1^2x_2^2x_3^2\le\tfrac13X_2q_4$, and since
$2p_0+1-\tfrac43X_2>0$ and $q_4\le\tfrac13X_2^2$,
\[
\begin{aligned}
 \frac{\mathcal X(p;x^2)}{X_2}&\ge18p_0(2p_0-1)-1-12(1-p_0)X_2+16X_2^2\\
 &=\left(4X_2-\frac32(1-p_0)\right)^2+\frac{135p_0^2-54p_0-13}{4}
 \ge\frac{16573}{78125}>0.
\end{aligned}
\]
On $r_0\in[\tfrac35,\tfrac{19}{25}]$, using $D_j^+\le D_{\max}$ and
$q_4\le\tfrac14(2r_0-1)^4$ gives
\[
 \frac{\mathcal X(p;x^2)}{X_2}\ge18p_0(2p_0-1)-1+36p_0X_2
   -\frac{9(2p_0+D_{\max})(2r_0-1)^4}{X_2},
\]
which is strictly increasing in $X_2>0$. On $[\tfrac35,\tfrac{16}{25}]$,
setting $D_{\max}=4p_0-1$ and
$X_2\ge X_{2,1}(r_0):=\tfrac1{15}(1-r_0)(15-11r_0)$ and multiplying by
$15X_{2,1}(r_0)>0$ gives the degree-six polynomial
$P_1(r_0)=15\bigl((36r_0^4-18r_0^2-1)X_{2,1}(r_0)+36r_0^2X_{2,1}(r_0)^2
-9(6r_0^2-1)(2r_0-1)^4\bigr)$, which is strictly concave on
$[\tfrac35,1]$ ($P_1''(\tfrac35+u)=-368208u^4-\tfrac{2057376}{5}u^3
-\tfrac{2990232}{25}u^2-\tfrac{1591704}{125}u-\tfrac{226078}{625}<0$ for all
$u\ge0$) with $P_1(\tfrac35)=\tfrac{3693}{78125}>0$ and
$P_1(\tfrac{16}{25})=\tfrac{216609912}{1220703125}>0$. On
$[\tfrac{16}{25},\tfrac{19}{25}]$, setting $D_{\max}=1$ and
$X_2\ge X_{2,2}(r_0):=2(1-r_0)(2r_0-1)$ and multiplying by
$\frac{X_{2,2}(r_0)}{2r_0-1}=2(1-r_0)>0$ gives the degree-five polynomial
$P_2(r_0)=72r_0^5-432r_0^4+432r_0^3-54r_0^2-52r_0+7$, which is strictly
concave on $[\tfrac35,1]$
($P_2''(\tfrac35+u)=-36\bigl(8u^2(9-5u)+\tfrac{288}{5}u+3\bigr)<0$ for
$u\in[0,\tfrac25]$) with $P_2(\tfrac{16}{25})=\tfrac{988047}{9765625}>0$
and $P_2(\tfrac{19}{25})=\tfrac{570453}{9765625}>0$. Thus
$\mathcal X(p;x^2)\ge0$ in all cases, proving \eqref{eq:fqe-critical-fidelity}.

Let $x=\sqrt{\Phi}$ and $c=(2+\sqrt2)/4$. The $32$ overlaps
$q_\eta=|\langle\eta,v\rangle|^2$ sum to $8$, since the frame operator is $8I$ (CO-1), so
$\Phi=\max_\eta q_\eta\ge1/4$ and $x\ge1/2$.
For $x\ge1/\sqrt2$, the first bound gives
$1-x^2\ge2c(1-x)$. For $1/2\le x\le1/\sqrt2$, the second gives
\[
 3/2-2x^2-2c(1-x)
 =2(1/\sqrt2-x)\bigl(x-(2-\sqrt2)/4\bigr)\ge0.
\]
Since $d^2=2(1-x)$, this proves \eqref{eq:fqe-critical-distance}.
Equation~\eqref{eq:fqe-orbit-surrogate} gives \eqref{eq:fqe-critical-entropy}, with
$\mathcal D=Q+5\tau/2$ on critical states and the bound just proved as its last step. Bell products have
$Q=\Phi=1/2$, hence $d^2=2-\sqrt2$; both constants are sharp.
The null-critical refinement is proved in the next paragraph.
\end{proof}

\paragraph{Two auxiliary Bloch spheres on the null-critical locus.}
Let $M=(m_1,m_2,m_3,m_4)$ be the unitary $4\times4$ matrix whose columns are the magic basis vectors of
Appendix~\ref{sec:fqe-normal-form} of the paper, written in the basis
$|00\rangle,|01\rangle,|10\rangle,|11\rangle$, so that $Mv=\sum_kv_km_{k+1}$ for $v\in\mathbb C^4$. For the critical
normal form $D_v$, regard $Mv$ as an auxiliary two-qubit vector.
Its coefficient matrix has determinant $\sum_jv_j^2/2$.
Thus $\tau=0$ is equivalent to $Mv=a\otimes b$ for unit auxiliary
qubits $a,b$. These are coordinates on the critical normal form, not
physical tensor factors of the original four-qubit state.
Let $r,s\in S^2$ be their Bloch vectors. Direct Bell probabilities give
\begin{equation}\label{eq:fqe-null-critical-bloch}
 Q=\frac12+\frac{27}{2}r_1r_2r_3s_1s_2s_3,\qquad
 \Phi=\frac14\max_{\substack{\varepsilon,\eta\in\{\pm1\}^3\\
                         \prod_j\varepsilon_j\eta_j=-1}}
 \left(1+\frac{\varepsilon\cdot r}{\sqrt3}\right)
 \left(1+\frac{\eta\cdot s}{\sqrt3}\right).
\end{equation}
Indeed, the four probabilities $|v_j|^2$ are
$(1+\sum_\alpha\chi_{j\alpha}r_\alpha s_\alpha)/4$, where
\[
 (\chi_{j\alpha})=
 \begin{pmatrix}1&-1&1\\1&1&-1\\-1&-1&-1\\-1&1&1\end{pmatrix}.
\]
Their second and third power sums in
\eqref{eq:fqe-global-intrinsic}, with $Z=D_v$ and $\operatorname{tr}(Z^TZ)=0$, give
the first identity. The 32 HS representatives become exactly the
auxiliary product rays whose Bloch vectors are cube vertices
$\varepsilon/\sqrt3,\eta/\sqrt3$ of opposite parity. The product-state
overlap formula and \eqref{eq:fqe-signed-diagonal-overlap} give the second
identity. These finite identities are checked independently by NC-1--NC-3.

Here is an elementary proof of the geometric claims in
\eqref{eq:fqe-null-critical-sharp}; NC-4 checks its scalar identities and constants.
Write $A=1-1/\sqrt3$, $B=1+1/\sqrt3$, and
$f_0=B^2/4=(2+\sqrt3)/6$.
First we prove $\Phi\ge f_0$. If the signs of the coordinates of
$r,s$ have opposite parity, align both cube vertices with those signs;
$\|r\|_1,\|s\|_1\ge1$ gives the claim. A zero coordinate allows either
parity. In the remaining case put $L=\|r\|_1$, $m=\min_j|r_j|$ and
define $f_r=1+L/\sqrt3$, $g_r=1+(L-2m)/\sqrt3$, with analogous
quantities for $s$. Flipping the sign at a smallest coordinate gives
\[
 4\Phi\ge\max(f_rg_s,g_rf_s)
                  \ge\sqrt{f_rg_rf_sg_s}.
\]
For the ordered coordinate moduli $x\ge y\ge z=m$ of either unit vector,
\[
 f_rg_r-B^2=\frac2{\sqrt3}(x+y-1)+\frac23(xy-z^2)\ge0.
\]
Here $xy\ge z^2$ and $(x+y)^2=1-z^2+2xy\ge1$.
This proves the fidelity floor in all cases.

If $r_1r_2r_3s_1s_2s_3\ge0$, then $Q\ge1/2$, and with $\Phi\ge f_0$ this gives
$Q\ge(1-\Phi)/(2(1-f_0))=\frac3{4-\sqrt3}(1-\Phi)$.
Otherwise put $\mu=3\sqrt3|r_1r_2r_3|$ and
$\nu=3\sqrt3|s_1s_2s_3|$, so $0\le\mu,\nu\le1$ and $2Q=1-\mu\nu$.
For a nonnegative unit triple $x,y,z$,
\[
 x+y+z-1
 =\frac{2(xy+yz+zx)}{x+y+z+1}
 \ge(9-3\sqrt3)xyz.
\]
For positive coordinates this follows from
$1/x+1/y+1/z\ge3\sqrt3$ and $x+y+z\le\sqrt3$;
zero coordinates follow by continuity.
Thus $\|r\|_1\ge1+(\sqrt3-1)\mu$ and similarly for $s$.
Aligning opposite-parity signs in \eqref{eq:fqe-null-critical-bloch} gives
$4\Phi\ge \Lambda:=(B+A\mu)(B+A\nu)$. The identity
\[
 \Lambda-B^2-(4-B^2)\mu\nu=AB\bigl(\mu(1-\nu)+\nu(1-\mu)\bigr)\ge0
\]
gives $1-\Phi\le(1-f_0)(1-\mu\nu)=2(1-f_0)Q$, the infidelity bound. In both cases the distance bound follows
pointwise: $1-\Phi=\frac{1+\sqrt\Phi}2\,d^2\ge\frac{1+\sqrt{f_0}}2\,d^2$, and
$\frac3{4-\sqrt3}\cdot\frac{1+\sqrt{f_0}}2=\frac{3+\sqrt3}4$ because $\sqrt{f_0}=B/2$.
The critical cluster representative $v=(1,i,0,0)/\sqrt2$ has
$Q=1/2$ and $\sqrt{\Phi}=B/2$, attaining all three geometric
bounds.


\begin{lemma}[Real-density geometry and sharp critical covering profile; exact certificate, RP, FP]
\label{lem:fqe-critical-real-density}
Write a unit critical normal vector as $v=x+iy$, with $x,y\in\mathbb R^4$,
and put $A=xx^T+yy^T$. Then $A$ ranges over exactly the real positive
semidefinite matrices of trace one and rank at most two, and
\begin{equation}\label{eq:fqe-critical-density-measurements}
 \lambda_{ji}=e_i^TH_jAH_j^Te_i,\qquad
 \tau=2\operatorname{tr}A^2-1.
\end{equation}
Here $H_0,H_1,H_2$ are the three real orthogonal cut transforms displayed
above. Their bases are mutually unbiased. For their twelve unit vectors
$h_{ji}=H_j^Te_i$ and rank-one projectors $B_{ji}=h_{ji}h_{ji}^T$,
\begin{equation}\label{eq:fqe-critical-real-tomography}
 A=\sum_{j,i}\lambda_{ji}B_{ji}-\tfrac12I,\qquad
 \sum_{j,i}\lambda_{ji}^2=\operatorname{tr}A^2+\tfrac12.
\end{equation}
For each coordinate axis $e_k$ and each of the four lines
$u=(0,\pm1,\pm1,\pm1)/\sqrt3$ in its orthogonal coordinate three-space,
put $P=e_ke_k^T+uu^T$ and $A_P=P/2$. There are sixteen such planes.
With $F_{\rm B}(A,B)=(\operatorname{tr}\sqrt{\sqrt B A\sqrt B})^2$,
\begin{equation}\label{eq:fqe-critical-bures-orbit}
 \Phi(v)=\max_P F_{\rm B}(A,A_P),\qquad
 F_{\rm B}(A,A_P)=\frac12\left[\operatorname{tr}(PA)
                    +2\sqrt{\det\bigl((e_k,u)^TA(e_k,u)\bigr)}\right].
\end{equation}
Thus $d(v)$ is exactly the minimum Bures distance from $A$ to these
sixteen real density matrices. The entire critical locus obeys the sharp bound
\begin{equation}\label{eq:fqe-critical-fidelity-floor}
 \Phi\ge f_{\rm cov}(\tau):=
       \frac13+\frac{\sqrt{1-\tau}}{2\sqrt3},\qquad
 \max_{\substack{\psi\ {\rm critical}\\\tau(\psi)=t}}d(\psi)^2
       =2-2\sqrt{f_{\rm cov}(t)}\quad(0\le t\le1).
\end{equation}
The bound is attained for every $t$ by the cluster-to-GHZ family below.
In particular every critical state has $d^2\le2-2/\sqrt3$, with equality for GHZ.
For any one such plane, put
$\lambda^P_{ji}=\operatorname{tr}(B_{ji}A_P)$ and
$D_P=\sum_{j,i}\lambda_{ji}\ln(\lambda_{ji}/\lambda^P_{ji})$.
The entropy gap has the exact form
\begin{equation}\label{eq:fqe-critical-relative-entropy}
 3\ln2\,\delta=D_P-\tfrac32\ln3\,J_P,\qquad
 J_P=1-\operatorname{tr}(PA)
    =\|A-A_P\|_F^2-\tau/2.
\end{equation}
\end{lemma}

\begin{proof}
The three physical cut vectors are $H_jv$, so their squared moduli are
the stated real measurements. Conversely, for any real $A$ of the
specified rank, its spectral decomposition gives
$v=\sqrt{a_1}e_1+i\sqrt{a_2}e_2$; rank one is included.
The identity $|v^Tv|^2=2\operatorname{tr}A^2-1$ proves the spin-flip claim.
The twelve real lines are the coordinate lines and the half-sign lines
of the regular $24$-cell. Their tight frame identity, valid for every
real symmetric matrix $C$, is
\[
 \sum_{j,i}\operatorname{tr}(B_{ji}C)B_{ji}
                    =C+\tfrac12\operatorname{tr}(C)I.
\]
Taking its trace pairing with $C$ proves both tomography formulas (RP-1).

The two HS rays associated with $P$ are $(e_k\pm iu)/\sqrt2$.
Their squared overlaps with $v$ have sum $\operatorname{tr}(PA)$ and
half their difference has squared value
$\det((e_k,u)^TA(e_k,u))$. Maximising over this pair therefore gives
\eqref{eq:fqe-critical-bures-orbit}, by the two-dimensional formula
$(\operatorname{tr}\sqrt C)^2=\operatorname{tr}C+2\sqrt{\det C}$.
The signed-Birkhoff reduction already proved shows that maximising
these sixteen pairwise maxima is the full physical HS fidelity.

For $q_\eta=|\langle\eta,v\rangle|^2$, the finite HS frame satisfies
\[
 \sum_{\eta\in\mathcal H}q_\eta=8,\qquad
 \sum_{\eta\in\mathcal H}q_\eta^2=\tfrac89(4-\tau).
\]
Both follow by expanding and cancelling odd sign powers (RP-3).
The stronger sharp profile \eqref{eq:fqe-critical-fidelity-floor} follows
from the area-weighted calibration proved below; the moment identity
alone does not give that profile.

Each reference distribution has one entry $1/2$ in each cut and nine
entries $1/6$ in total. The sum of the three projectors selecting its
large entries is $3P/2$ (RP-3). Expanding $D_P$ then gives
\eqref{eq:fqe-critical-relative-entropy}. Finally
$\operatorname{tr}A_P^2=1/2$ and
$2\operatorname{tr}(AA_P)=\operatorname{tr}(AP)$ give the leakage formula.
All matrix and compression identities here are independently checked by RP;
the covering calibration uses FP.
\end{proof}

\paragraph{Area-weighted covering calibration.}
We prove the profile in \eqref{eq:fqe-critical-fidelity-floor} using an
explicit probability measure on the sixteen HS planes. Its only
polynomial positivity input is the following finite exact certificate.

\begin{lemma}[Four-point cubic inequality; exact rational certificate]
\label{lem:fqe-four-point-cubic}
If $X,Y\in\mathbb R^4$ satisfy
$\sum_iX_i=\sum_iY_i=0$ and $\sum_iX_i^2=\sum_iY_i^2=4$, then
\begin{equation}\label{eq:fqe-four-point-cubic}
 \sum_{i,j=1}^4|X_i-Y_j|^3\ge64.
\end{equation}
Equality holds at $X=Y=(-1,-1,1,1)$.
\end{lemma}

\begin{proof}
Put $X=Er$, $Y=Es$, where
\[
 E=\begin{pmatrix}1&1&1\\1&-1&-1\\-1&1&-1\\-1&-1&1\end{pmatrix},
 \qquad E^TE=4I,
\]
so that $\|r\|=\|s\|=1$. Every zero-sum four-vector is represented.
Sort $X$ and $Y$ separately and interleave their entries in increasing
order. There are $\binom84=70$ words with four letters of each kind.
Swapping $X,Y$ and simultaneously negating both vectors reduce these
to exactly twenty-three classes (FP-1). Ties may use either order.
For each word let $z_1,\ldots,z_8$ be its ordered entries and
$g_i=z_{i+1}-z_i\ge0$. The absolute values then give a cubic polynomial
$p=\sum_{ij}|X_i-Y_j|^3/16-4$.
For each of the twenty-three words the supplied rational data give
the identity
\begin{equation}\label{eq:fqe-four-point-cubic-sos}
 p=B_0^TG_0B_0+\sum_{i=1}^7g_i B_i^TG_iB_i
       +(\|r\|^2-1)\ell_r+(\|s\|^2-1)\ell_s.
\end{equation}
The $\ell$ are quadratic polynomials and all eight $G_i$ are positive
definite. FP-2 rebuilds all $210$ coefficients through degree four
in each class and proves positivity of all $184$ Grams using integer
arithmetic. Thus the right side is nonnegative on each ordering cone
and the two spheres, proving the full inequality. The bases and exact
matrix representation are specified in Section~\ref{sec:sup-verifiers}.
The stated equality follows by substitution.
\end{proof}

Let $U=(U_1,U_2)$ be a real orthonormal $4\times2$ frame with plane projector
$\Pi=UU^T$. For each HS plane $P$ use its orthonormal frame
$V_P=(e_k,u)$ and put $c_P=|\det(U^TV_P)|$. The sixteen planes obey
\begin{equation}\label{eq:fqe-area-weighted-fusion}
 \sum_Pc_P^2=\frac83,\qquad
 \sum_Pc_P^2P=\frac89(I+\Pi).
\end{equation}
Indeed, the bivector $\Delta=U_1U_2^T-U_2U_1^T$ satisfies $\Delta\Delta^T=\Delta^T\Delta=\Pi$, and
$\det(U^TV_P)=e_k^T\Delta u$. Fix $k$, put $g=\Delta^Te_k$ (so $g_k=0$,
$|g|^2=\Pi_{kk}$) and sum over the four lines $u$ for this $k$, using $u_j=\pm1/\sqrt3$ ($j\ne k$):
$\sum_u(g\cdot u)^2=\frac43|g|^2$ and
$\sum_u(g\cdot u)^2uu^T=\frac49|g|^2(I-e_ke_k^T)+\frac89(gg^T-\operatorname{diag}(g_j^2))$.
Summing over $k$ with $\sum_kgg^T=\Pi$ and $\sum_k\operatorname{diag}(g_j^2)=\operatorname{diag}(\Pi)$ gives
both identities. FP-3 also checks them without
assuming orthonormality first, using the homogeneous squared area and
its associated matrix. Therefore the explicit weights
$w_P=3c_P^2/8$ satisfy
\begin{equation}\label{eq:fqe-plane-calibration}
 \sum_Pw_P=1,\qquad \sum_Pw_PU^TPU=\frac23I_2,\qquad
 \sum_Pw_Pc_P\ge\frac1{\sqrt3}.
\end{equation}

To prove the last inequality, let $M$ be the magic-basis matrix defined before
\eqref{eq:fqe-null-critical-bloch}. Form the null critical vector
$z=(U_1+iU_2)/\sqrt2$. Its auxiliary product representation
$Mz=a\otimes b$ has unit Bloch vectors $r,s$, as in
\eqref{eq:fqe-null-critical-bloch}.
For the two HS rays $\eta_P=(e_k+iu)/\sqrt2$ and $\bar\eta_P$,
expansion gives
\[
 c_P=\bigl||\langle\eta_P,z\rangle|^2
                 -|\langle\bar\eta_P,z\rangle|^2\bigr|.
\]
In the cube-pair description of the HS rays given after \eqref{eq:fqe-null-critical-bloch}, conjugation reverses both cube vertices.
Indeed $M=\epsilon^{\otimes2}\bar M$, where
$\epsilon=\left(\begin{smallmatrix}0&1\\-1&0\end{smallmatrix}\right)$
(FP-3), and the qubit spin flip reverses the Bloch vector.
Choose the four even-parity cube vertices $e_i$, the rows of $E$.
The sixteen conjugate pairs are represented once by
$(\varepsilon,\eta)=(e_i,-e_j)$, and hence
\[
 c_P=\frac{|e_i\cdot r-e_j\cdot s|}{2\sqrt3}.
\]
The four numbers $X_i=e_i\cdot r$, and likewise $Y_j=e_j\cdot s$,
have sum zero and sum of squares four. Lemma~\ref{lem:fqe-four-point-cubic}
therefore gives
\[
 \sum_Pc_P^3=\frac{\sum_{ij}|X_i-Y_j|^3}{24\sqrt3}
       \ge\frac8{3\sqrt3},
\]
which proves the last calibration inequality.

Now write the real density matrix of any critical state as
$A=U\Lambda U^T$, with $\operatorname{tr}\Lambda=1$ and
$\det\Lambda=(1-\tau)/4$. Rank one is included by completing its
support to any plane. The exact Bures formula gives
\begin{align*}
 \Phi
 &\ge\sum_Pw_P F_{\rm B}(A,P/2)\\
 &=\frac12\operatorname{tr}\!\left(\Lambda\sum_Pw_PU^TPU\right)
                  +\sqrt{\det\Lambda}\sum_Pw_Pc_P\\
 &\ge\frac13+\frac{\sqrt{1-\tau}}{2\sqrt3}.
\end{align*}
This proves the profile for the entire critical locus.

Finally, for $a,b\ge0$, $a+b>0$, the critical normal family
\begin{equation}\label{eq:fqe-covering-equality-family}
 v=\frac{(a+ib,-a+ib,0,0)}{\sqrt{2(a^2+b^2)}},\qquad
 \tau=\frac{(a^2-b^2)^2}{(a^2+b^2)^2},\qquad
 \Phi=\frac13+\frac{ab}{\sqrt3(a^2+b^2)}
\end{equation}
attains equality. FP-4 checks every one of the $32$ ray overlaps against
this maximum. The parameters cover $0\le\tau\le1$, proving pointwise
optimality and the conditional covering radius. The endpoint $a=b$
is a cluster state; $b=0$ is GHZ up to LU.


This model expresses both the entropy and the optimised orbit distance exactly.
In particular, a bound on the classical relative entropy $D_P$ that
controls the explicit plane leakage $J_P$ can improve the entropy
constant without first improving a moment surrogate. Its domain is a
real rank-two density matrix; replacing it by all real density matrices
would admit $I/4$, whose entropy exceeds the critical-state maximum.
